\section{The Emptiness of Meaning}

\begin{quotex}
What is life? Its opposite is petrification, the immobile duration of the identical. Therefore, for life to exist, it is necessary to have a sequential change of non-identical states.
\flright{\textsc{Pavel Florensky}}
\end{quotex}


In \emph{The Meaning of Idealism}, \textbf{Pavel Florensky} develops some useful insights based on Porphyry's scheme of the spheres of being. There are four of them:

\begin{enumerate}


\item \textbf{Ideas}: The universal concepts that express the One and the Many.

\item \textbf{Things}: The appearance of ideas in the world.

\item \textbf{Concepts}: The knowability of things.

\item \textbf{Terms}: The expressibility of the concepts in language.

\end{enumerate}


Differing philosophical solutions to the problem of universals can be described with this schema. For example, Platonic realism accepts all four spheres while total nihilism rejects them all. Aristotelianism admits 2, 3, and 4; the Ideas exist in the Thing, not independently.

In our day, Platonic realism is pejoratively dismissed as ``essentialism''. The dominant doctrine is nominalism, which denies (1) universal ideas, (2) their appearance in the world, and therefore (3) their knowability. Rather, what we know as ideas are the (4) ``words'' or terms that we use to describe things and events; they have no existences beyond our sense experience. Hence, knowledge is merely conventional --- i.e., general consensus --- not real.

That is why there is so much emphasis over words, why there is social pressure to use the ``correct'' word (not the most real word, but the one that is socially acceptable). For nominalism, only power is real and the control of words is tantamount to the control of the ideas that may be expressed.

Since materiality is the principle of multiplicity, nominalism renames it ``diversity''. Since it cannot admit the existence of the transcendent One, the principle of Unity is the general agreement to use the same ``words'' as descriptors of reality. That is why nominalists all sound alike, even repeating the same phrases over and over. That is also why their method is coercion rather than rational arguments.

This represents the end of a civilization. Everything is reduced to matter, dead matter.\footnote{This is the prolegomenon to a full review of the book. [It seems that such review was never done.]}


\flrightit{Posted on 2022-03-26 by Cologero}
